
ewcommand\QCMtexteCorr[7]{#1}
{% load latex %}

\documentclass[12pt,a4paper]{article}

% ===== Encodage & langue =====
\usepackage[utf8]{inputenc}
\usepackage[T1]{fontenc}
\usepackage[french]{babel}

% ===== Mise en page & graphismes =====
\usepackage{geometry}
\geometry{top=1.2cm, bottom=2cm, left=2cm, right=2cm}
\usepackage{graphicx}
\usepackage{array} % pour >{\raggedleft\arraybackslash} dans tabular
\usepackage{tikz}

% ===== Divers =====
\usepackage{hyperref}
\usepackage{amsmath,amssymb}
\usepackage{xcolor}
% ===== CHEMISTRY & ROBUSTNESS =====
% Compteur global des exercices
\newcounter{exoTop}
\makeatletter
% Charger mhchem si pas déjà chargé
\@ifpackageloaded{mhchem}{}{%
  \IfFileExists{mhchem.sty}{%
    \usepackage[version=4]{mhchem}%
    \typeout{DEBUG: mhchem.sty FOUND -> LOADED}%
  }{%
    \typeout{DEBUG: mhchem.sty NOT FOUND -> FALLBACK ACTIVE}%
  }%
}
% Fallback minimal si, pour une raison quelconque, mhchem n'est pas (ou n'a pas pu être) chargé :
\@ifpackageloaded{mhchem}{%
  \typeout{DEBUG: mhchem package is loaded at preamble.}%
}{%
  \providecommand{\ce}[1]{\textrm{#1}}% rend lisible sans planter
}
\makeatother

% Atténue les Overfull \hbox (non bloquant)
\setlength{\emergencystretch}{2em}

\usepackage[none]{hyphenat}
\sloppy

% --- Macro standard pour ligne de réponse (flashcards) ---
\newcommand{\AnswerLine}[1]{%
  \par\smallskip\noindent\textbf{Réponse : }#1\par%
}

% ===== Unicode fixes (indices / exposants / symboles courants) =====
\usepackage{newunicodechar}
% Mapping Unicode -> LaTeX (grec, symboles courants)
\DeclareUnicodeCharacter{03BB}{\ensuremath{\lambda}} % λ
\DeclareUnicodeCharacter{03BC}{\ensuremath{\mu}}     % μ
\DeclareUnicodeCharacter{03C0}{\ensuremath{\pi}}     % π
\DeclareUnicodeCharacter{03B1}{\ensuremath{\alpha}}  % α
\DeclareUnicodeCharacter{03B2}{\ensuremath{\beta}}   % β
\DeclareUnicodeCharacter{03B3}{\ensuremath{\gamma}}  % γ
\DeclareUnicodeCharacter{03B4}{\ensuremath{\delta}}  % δ
\DeclareUnicodeCharacter{03B8}{\ensuremath{\theta}}  % θ
\DeclareUnicodeCharacter{03C3}{\ensuremath{\sigma}}  % σ
\DeclareUnicodeCharacter{03C4}{\ensuremath{\tau}}    % τ
\DeclareUnicodeCharacter{03C6}{\ensuremath{\phi}}    % φ
\DeclareUnicodeCharacter{03C9}{\ensuremath{\omega}}  % ω
\DeclareUnicodeCharacter{0394}{\ensuremath{\Delta}}  % Δ
\DeclareUnicodeCharacter{03A3}{\ensuremath{\Sigma}}  % Σ
\DeclareUnicodeCharacter{03A9}{\ensuremath{\Omega}}  % Ω
\DeclareUnicodeCharacter{00B0}{\ensuremath{^\circ}}  % °

% Indices (₀..₉)
\newunicodechar{₀}{$_0$}
\newunicodechar{₁}{$_1$}
\newunicodechar{₂}{$_2$}
\newunicodechar{₃}{$_3$}
\newunicodechar{₄}{$_4$}
\newunicodechar{₅}{$_5$}
\newunicodechar{₆}{$_6$}
\newunicodechar{₇}{$_7$}
\newunicodechar{₈}{$_8$}
\newunicodechar{₉}{$_9$}

% Exposants (⁰ ¹ ² ³ ⁴..⁹)
\newunicodechar{⁰}{$^0$}
\newunicodechar{¹}{$^1$}
\newunicodechar{²}{$^2$}
\newunicodechar{³}{$^3$}
\newunicodechar{⁴}{$^4$}
\newunicodechar{⁵}{$^5$}
\newunicodechar{⁶}{$^6$}
\newunicodechar{⁷}{$^7$}
\newunicodechar{⁸}{$^8$}
\newunicodechar{⁹}{$^9$}

% Autres utiles
\newunicodechar{°}{\textdegree}
\newunicodechar{µ}{$\mu$}
\newunicodechar{×}{$\times$}
\newunicodechar{−}{$-$}  % signe moins U+221

% =========================================================
% Variables côté Jinja2 -> macros LaTeX attendues par l'entête
% =========================================================
% type ds / serie
% Type DS (booléen LaTeX)
\newif\ifDS
{% if type == "ds" %}\DStrue{% else %}\DSfalse{% endif %}

% titre / barème / sous-titres
\def\titre{ {{ titre|tex }} }
\def\baremeGlobal{ {{ barreme_obtenu|default:barreme_cible|default:0 }} }
\def\sousTitreUn{ {{ sous_titre_1|tex }} }
\def\sousTitreDeux{ {{ sous_titre_2|tex }} }

% --- Logo: flag + chemin (pas d'include si vide) ---
\newif\ifHAVELOGO
{% if header_logo_path and header_logo_path != "" %}
  \HAVELOGOtrue
{% else %}
  \HAVELOGOfalse
{% endif %}



% Définition du chemin du logo pour header.tex
\def\headerLogoPath{ {{ header_logo_path|default:"" }} }

% DEBUG: Afficher le chemin du logo
\typeout{DEBUG: header_logo_path = {{ header_logo_path }}}
\typeout{DEBUG: headerLogoPath = \headerLogoPath}

% lignes du pavé gauche si pas de logo
\def\headerLineUn{ {{ header_line1|default:""|tex }} }
\def\headerLineDeux{ {{ header_line2|default:""|tex }} }
\def\headerLineTrois{ {{ header_line3|default:""|tex }} }

% classe / nom (tu peux les alimenter via payload si tu veux)
\def\classe{ {{ classe|default:""|tex }} }
\def\nomPrenom{ Nom et Prénom : }

% énoncé (0) ou correction (1)
\def\isCorrection{ {{ isCorrection|default:0 }} }

\begin{document}

% ===== ENTÊTE =====
% generator/templates/generator/pdf/header.tex
% Attend que main.tex ait défini :
% \typeDS, \ds, \serie, \titre, \baremeGlobal, \sousTitreUn, \sousTitreDeux,
% \headerLineUn, \headerLineDeux, \headerLineTrois, \classe, \nomPrenom, \isCorrection
% ainsi que \ifHAVELOGO et la variable de template {{ header_logo_path }}

\noindent
\begin{tabular*}{\textwidth}{@{\extracolsep{\fill}} p{0.68\textwidth} p{0.28\textwidth}}
  \normalsize \nomPrenom & \raggedleft \normalsize Classe : \classe\\
\end{tabular*}

\vspace{0.25cm}

\noindent
\begin{tikzpicture}[x=1cm,y=1cm, line width=0.5pt]

  % Dimensions
  \pgfmathsetmacro{\X}{\textwidth/1cm}
  \def\H{2}

  % Cadres
  \draw (0,0) rectangle (\X,\H);     % grand rectangle
  \draw (0,0) rectangle (2,\H);      % pavé gauche

  % Pavé droit (barème) uniquement en DS
  \ifx\typeDS\ds
    \draw (\X-2,0) rectangle (\X,\H);
    \draw (\X-2,0) -- (\X,\H);
  \fi

  % Contenu pavé gauche : logo OU 3 lignes
  \ifHAVELOGO
    \node at (1,1) {\includegraphics[width=2cm,height=2cm,keepaspectratio]{\headerLogoPath}};
  \else
    \node[align=center] at (1,1.35) {\large\bfseries \headerLineUn};
    \node[align=center] at (1,1.00) {\small \headerLineDeux};
    \node[align=center] at (1,0.65) {\small \headerLineTrois};
  \fi

  % Abscisse du bloc central (diffère entre DS et Série)
  \ifx\typeDS\ds
    \pgfmathsetmacro{\XC}{\X/2}
  \else
    \pgfmathsetmacro{\XC}{(\X+2)/2}
  \fi

  % Titre et sous-titres
  \node[align=center] at (\XC,1.45) {\Large \bfseries \titre \ifnum\isCorrection=1 ~|~ Correction\fi};
  \node[align=center] at (\XC,0.95) {\normalsize \sousTitreUn};
  \node[align=center] at (\XC,0.55) {\normalsize \sousTitreDeux};

  % Barème (DS)
  \ifx\typeDS\ds
    \node[align=center] at (\X-0.5,0.7) {\Large \bfseries \baremeGlobal};
  \fi

\end{tikzpicture}


% ===== INFORMATION BARÈME AJUSTÉ =====
{% if barreme_obtenu != barreme_cible %}
\vspace{0.5cm}
\begin{center}
\fbox{\parbox{0.8\textwidth}{
\centering
\textbf{Note :} Le barème a été ajusté de {{ barreme_cible }} points à {{ barreme_obtenu }} points\\
en raison du nombre limité d'exercices disponibles.
}}
\end{center}
\vspace{0.5cm}
{% endif %}

% ===== SECTION A : QCM =====
{% if qcms|length > 0 %}
{% load latex %}

% ====== QUESTIONS DE COURS : EN-TÊTE ======
{% if has_qcms or has_flashcards %}
  % Démarre l'exercice 1 une seule fois ici (si QCM présents)
  \setcounter{exoTop}{1}
  \noindent\textbf{Exercice \Roman{exoTop} -- Questions de cours}\par\vspace{4pt}
{% endif %}

% ====== TITRE DE SOUS-SECTION (A. si QCM + Flashcards, sinon sans lettre) ======
{% if has_qcms %}
  \noindent\textbf{
    {% if has_flashcards %}A. {% endif %}QCM
    {% if qcm_points_sum %} ({{ qcm_points_sum|floatformat:0 }} pt){% endif %}
  }\par\vspace{6pt}
{% endif %}

% ============================
% COMPTEUR
% ============================
\newcounter{qcmno}

% ============================
% MACROS ENONCE
% ============================

% QCM texte pur
\newcommand{\QCMtexte}[6]{%
  \stepcounter{qcmno}%
  \noindent\textbf{\theqcmno.}~#1

  \vspace{0.5em}
  \begin{tabular}{@{}p{0.45\linewidth}@{}p{0.45\linewidth}@{}>{\raggedleft\arraybackslash}p{0.10\linewidth}@{}}
    A) #2 & B) #3 & \\
    \multicolumn{2}{c}{} & \textbf{#6 pt} \\
    C) #4 & D) #5 & \\
  \end{tabular}
  \vspace{1.2em}
}

% QCM texte + image
\newcommand{\QCMimage}[7]{%
  \stepcounter{qcmno}%
  \noindent
  \begin{minipage}[t]{0.69\linewidth}
    \textbf{\theqcmno.}~#1\par
  \end{minipage}%
  \hspace{0.02\linewidth}%
  \begin{minipage}[t]{0.29\linewidth}
    \vspace{-2pt}%
    \includegraphics[width=\linewidth]{#2}
  \end{minipage}

  \vspace{0.5em}
  \begin{tabular}{@{}p{0.45\linewidth}@{}p{0.45\linewidth}@{}>{\raggedleft\arraybackslash}p{0.10\linewidth}@{}}
    A) #3 & B) #4 & \\
    \multicolumn{2}{c}{} & \textbf{#7 pt} \\
    C) #5 & D) #6 & \\
  \end{tabular}
  \vspace{1.2em}
}

% ============================
% MACROS CORRECTION
% ============================

% QCM texte pur (bonne réponse en gras)
\newcommand{\QCMtexteCorr}[7]{%
  \stepcounter{qcmno}%
  \noindent\textbf{\theqcmno.}~#1

  \vspace{0.5em}
  \begin{tabular}{@{}p{0.45\linewidth}@{}p{0.45\linewidth}@{}>{\raggedleft\arraybackslash}p{0.10\linewidth}@{}}
    A) \ifnum#7=1 \textbf{#2} \else #2 \fi &
    B) \ifnum#7=2 \textbf{#3} \else #3 \fi & \\
    \multicolumn{2}{c}{} & \textbf{#6 pt} \\
    C) \ifnum#7=3 \textbf{#4} \else #4 \fi &
    D) \ifnum#7=4 \textbf{#5} \else #5 \fi & \\
  \end{tabular}
  \vspace{1.2em}
}

% QCM image (bonne réponse en gras)
\newcommand{\QCMimageCorr}[8]{%
  \stepcounter{qcmno}%
  \noindent
  \begin{minipage}[t]{0.69\linewidth}
    \textbf{\theqcmno.}~#1\par
  \end{minipage}%
  \hspace{0.02\linewidth}%
  \begin{minipage}[t]{0.29\linewidth}
    \vspace{-2pt}%
    \includegraphics[width=\linewidth]{#2}
  \end{minipage}

  \vspace{0.5em}
  \begin{tabular}{@{}p{0.45\linewidth}@{}p{0.45\linewidth}@{}>{\raggedleft\arraybackslash}p{0.10\linewidth}@{}}
    A) \ifnum#8=1 \textbf{#3} \else #3 \fi &
    B) \ifnum#8=2 \textbf{#4} \else #4 \fi & \\
    \multicolumn{2}{c}{} & \textbf{#7 pt} \\
    C) \ifnum#8=3 \textbf{#5} \else #5 \fi &
    D) \ifnum#8=4 \textbf{#6} \else #6 \fi & \\
  \end{tabular}
  \vspace{1.2em}
}

% ============================
% SECTION
% ============================

{% for item in qcms %}
  {% if item.has_image %}
    {% if is_correction %}
\QCMimageCorr{ {{ item.statement|tex }} }
{ {{ item.image_path }} }
{ {{ item.A|tex }} }
{ {{ item.B|tex }} }
{ {{ item.C|tex }} }
{ {{ item.D|tex }} }
{ {{ item.points|floatformat:0 }} }
{ {{ item.correct_index }} }
    {% else %}
\QCMimage{ {{ item.statement|tex }} }
{ {{ item.image_path }} }
{ {{ item.A|tex }} }
{ {{ item.B|tex }} }
{ {{ item.C|tex }} }
{ {{ item.D|tex }} }
{ {{ item.points|floatformat:0 }} }
    {% endif %}
  {% else %}
    {% if is_correction %}
\QCMtexteCorr{ {{ item.statement|tex }} }
{ {{ item.A|tex }} }
{ {{ item.B|tex }} }
{ {{ item.C|tex }} }
{ {{ item.D|tex }} }
{ {{ item.points|floatformat:0 }} }
{ {{ item.correct_index }} }
    {% else %}
\QCMtexte{ {{ item.statement|tex }} }
{ {{ item.A|tex }} }
{ {{ item.B|tex }} }
{ {{ item.C|tex }} }
{ {{ item.D|tex }} }
{ {{ item.points|floatformat:0 }} }
    {% endif %}
  {% endif %}
{% empty %}
\noindent\emph{Aucun QCM disponible.}
{% endfor %}

{% endif %}

% ===== SECTION B : Flashcards =====
{% if flashcards|length > 0 %}
{% load latex %}

{% if has_flashcards %}
  % Si pas de QCM, on imprime l'entête "Exercice 1 -- Questions de cours" ici
  {% if not has_qcms %}
    \setcounter{exoTop}{1}
    \noindent\textbf{Exercice \Roman{exoTop} -- Questions de cours}\par\vspace{4pt}
  {% endif %}

  % Titre de sous-section : "B." seulement s'il y a AUSSI des QCM
  \noindent\textbf{
    {% if has_qcms %}B. {% endif %}Questions
    {% if flashcards_points_sum %} ({{ flashcards_points_sum|floatformat:0 }} pt){% endif %}
  }\par\vspace{6pt}

  % Liste des cartes
  {% for card in flashcards %}
    \vspace{0.5em}
    \noindent
    \begin{tabular}{@{}p{0.90\linewidth}@{}>{\raggedleft\arraybackslash}p{0.10\linewidth}@{}}
      \textbf{ {{ forloop.counter }}.\ }{{ card.question_html|safe|html2tex }} &
      {% if card.points %}\textbf{ {{ card.points|floatformat:0 }} pt}{% endif %} \\
    \end{tabular}

    {% if is_correction and card.answer_html %}
      \par\smallskip\noindent\textcolor{green}{\textbf{Réponse : }} {{ card.answer_html|safe|html2tex }}\par
    {% endif %}
    \vspace{1.2em}
  {% empty %}
    \noindent\emph{Aucune flashcard sélectionnée.}
  {% endfor %}
{% endif %}

{% endif %}

% ===== SECTION C : Exercices =====
{% if exercices|length > 0 %}
{% load latex %}{% load extras %}{% load param_tags %}

% Compteur des sous-questions par exercice
\newcounter{exQno}

% Pas de \section* ici : on enchaîne après "Questions de cours"

{% for ex in exercices %}
  % Passe à l'exercice suivant (2, 3, …)
  \addtocounter{exoTop}{1}
  \setcounter{exQno}{0}

  % Ligne de titre alignée comme QCM (90% titre | 10% barème)
  \noindent
  \begin{minipage}[t]{0.90\linewidth}
    \textbf{Exercice \Roman{exoTop} | {{ ex.title|html2tex }}}
  \end{minipage}%
  \begin{minipage}[t]{0.10\linewidth}
    \raggedleft
    {% if ex.total_points %}\textbf{ ({{ ex.total_points|floatformat:0 }} pt)}{% endif %}
  \end{minipage}
  \par\vspace{6pt}

  {% if ex.rendered_blocks %}
    {% for block in ex.rendered_blocks %}
      {% if block.block_type == 'enonce' or block.block_type == 's100' %}
        \noindent {{ block.rendered_html|html2tex }}\vspace{8pt}

      {% elif block.block_type == 's70_30' %}
        \noindent
        \begin{minipage}[t]{0.70\linewidth}
          {{ block.rendered_html|html2tex_no_images }}
        \end{minipage}%
        \hspace{0.02\linewidth}%
        \begin{minipage}[t]{0.28\linewidth}
          \vspace{-2pt}\centering
          {% if ex.image_path %}\includegraphics[width=\linewidth,height=0.28\textheight,keepaspectratio]{ {{ ex.image_path }} }{% endif %}
        \end{minipage}
        \vspace{8pt}

      {% elif block.block_type == 's50_50' %}
        \noindent
        \begin{minipage}[t]{0.48\linewidth}
          {{ block.rendered_html|html2tex_no_images }}
        \end{minipage}%
        \hspace{0.04\linewidth}%
        \begin{minipage}[t]{0.48\linewidth}
          \vspace{-2pt}\centering
          {% if ex.image_path %}\includegraphics[width=\linewidth,height=0.48\textheight,keepaspectratio]{ {{ ex.image_path }} }{% endif %}
        \end{minipage}
        \vspace{8pt}

      {% elif block.block_type == 's75_25' %}
        \noindent
        \begin{minipage}[t]{0.75\linewidth}
          {{ block.rendered_html|html2tex_no_images }}
        \end{minipage}%
        \hspace{0.01\linewidth}%
        \begin{minipage}[t]{0.24\linewidth}
          \vspace{-2pt}\centering
          {% if ex.image_path %}\includegraphics[width=\linewidth,height=0.24\textheight,keepaspectratio]{ {{ ex.image_path }} }{% endif %}
        \end{minipage}
        \vspace{8pt}

      {% elif block.block_type == 'question' %}
        \refstepcounter{exQno}%
        \vspace{0.5em}
        \noindent
        \begin{minipage}[t]{0.90\linewidth}
          \textbf{\theexQno. }\ignorespaces
          {% for child in block.value.content %}{% render_with_params child ex.param_values as qhtml %}{{ qhtml|strip_badges|html2tex }}{% endfor %}
        \end{minipage}%
        \hfill
        \begin{minipage}[t]{0.10\linewidth}
          \raggedleft
          {% if block.value.points %}\textbf{ {{ block.value.points|floatformat:0 }} pt}{% endif %}
        \end{minipage}
        \par
        \vspace{1.2em}

        {% if is_correction and block.value.solution %}
          \noindent\textbf{Solution :}
          \par\vspace{4pt}
          {% for child in block.value.solution %}{% render_with_params child ex.param_values as solhtml %}
            \noindent {{ solhtml|html2tex }}\par
          {% endfor %}
          \vspace{8pt}
        {% endif %}

      {% else %}
        {{ block.rendered_html|html2tex }}\vspace{4pt}
      {% endif %}
    {% endfor %}
  {% endif %}

  \vspace{12pt}
{% empty %}
  \noindent\emph{Aucun exercice sélectionné.}
{% endfor %}

{% endif %}

\end{document}
% generator/templates/generator/pdf/header.tex
% Attend que main.tex ait défini :
% \typeDS, \ds, \serie, \titre, \baremeGlobal, \sousTitreUn, \sousTitreDeux,
% \headerLineUn, \headerLineDeux, \headerLineTrois, \classe, \nomPrenom, \isCorrection
% ainsi que \ifHAVELOGO et la variable de template {{ header_logo_path }}

\noindent
\begin{tabular*}{\textwidth}{@{\extracolsep{\fill}} p{0.68\textwidth} p{0.28\textwidth}}
  \normalsize \nomPrenom & \raggedleft \normalsize Classe : \classe\\
\end{tabular*}

\vspace{0.25cm}

\noindent
\begin{tikzpicture}[x=1cm,y=1cm, line width=0.5pt]

  % Dimensions
  \pgfmathsetmacro{\X}{\textwidth/1cm}
  \def\H{2}

  % Cadres
  \draw (0,0) rectangle (\X,\H);     % grand rectangle
  \draw (0,0) rectangle (2,\H);      % pavé gauche

  % Pavé droit (barème) uniquement en DS
  \ifx\typeDS\ds
    \draw (\X-2,0) rectangle (\X,\H);
    \draw (\X-2,0) -- (\X,\H);
  \fi

  % Contenu pavé gauche : logo OU 3 lignes
  \ifHAVELOGO
    \node at (1,1) {\includegraphics[width=2cm,height=2cm,keepaspectratio]{\headerLogoPath}};
  \else
    \node[align=center] at (1,1.35) {\large\bfseries \headerLineUn};
    \node[align=center] at (1,1.00) {\small \headerLineDeux};
    \node[align=center] at (1,0.65) {\small \headerLineTrois};
  \fi

  % Abscisse du bloc central (diffère entre DS et Série)
  \ifx\typeDS\ds
    \pgfmathsetmacro{\XC}{\X/2}
  \else
    \pgfmathsetmacro{\XC}{(\X+2)/2}
  \fi

  % Titre et sous-titres
  \node[align=center] at (\XC,1.45) {\Large \bfseries \titre \ifnum\isCorrection=1 ~|~ Correction\fi};
  \node[align=center] at (\XC,0.95) {\normalsize \sousTitreUn};
  \node[align=center] at (\XC,0.55) {\normalsize \sousTitreDeux};

  % Barème (DS)
  \ifx\typeDS\ds
    \node[align=center] at (\X-0.5,0.7) {\Large \bfseries \baremeGlobal};
  \fi

\end{tikzpicture}
{% load latex %}

% ====== QUESTIONS DE COURS : EN-TÊTE ======
{% if has_qcms or has_flashcards %}
  % Démarre l'exercice 1 une seule fois ici (si QCM présents)
  \setcounter{exoTop}{1}
  \noindent\textbf{Exercice \Roman{exoTop} -- Questions de cours}\par\vspace{4pt}
{% endif %}

% ====== TITRE DE SOUS-SECTION (A. si QCM + Flashcards, sinon sans lettre) ======
{% if has_qcms %}
  \noindent\textbf{
    {% if has_flashcards %}A. {% endif %}QCM
    {% if qcm_points_sum %} ({{ qcm_points_sum|floatformat:0 }} pt){% endif %}
  }\par\vspace{6pt}
{% endif %}

% ============================
% COMPTEUR
% ============================
\newcounter{qcmno}

% ============================
% MACROS ENONCE
% ============================

% QCM texte pur
\newcommand{\QCMtexte}[6]{%
  \stepcounter{qcmno}%
  \noindent\textbf{\theqcmno.}~#1

  \vspace{0.5em}
  \begin{tabular}{@{}p{0.45\linewidth}@{}p{0.45\linewidth}@{}>{\raggedleft\arraybackslash}p{0.10\linewidth}@{}}
    A) #2 & B) #3 & \\
    \multicolumn{2}{c}{} & \textbf{#6 pt} \\
    C) #4 & D) #5 & \\
  \end{tabular}
  \vspace{1.2em}
}

% QCM texte + image
\newcommand{\QCMimage}[7]{%
  \stepcounter{qcmno}%
  \noindent
  \begin{minipage}[t]{0.69\linewidth}
    \textbf{\theqcmno.}~#1\par
  \end{minipage}%
  \hspace{0.02\linewidth}%
  \begin{minipage}[t]{0.29\linewidth}
    \vspace{-2pt}%
    \includegraphics[width=\linewidth]{#2}
  \end{minipage}

  \vspace{0.5em}
  \begin{tabular}{@{}p{0.45\linewidth}@{}p{0.45\linewidth}@{}>{\raggedleft\arraybackslash}p{0.10\linewidth}@{}}
    A) #3 & B) #4 & \\
    \multicolumn{2}{c}{} & \textbf{#7 pt} \\
    C) #5 & D) #6 & \\
  \end{tabular}
  \vspace{1.2em}
}

% ============================
% MACROS CORRECTION
% ============================

% QCM texte pur (bonne réponse en gras)
\newcommand{\QCMtexteCorr}[7]{%
  \stepcounter{qcmno}%
  \noindent\textbf{\theqcmno.}~#1

  \vspace{0.5em}
  \begin{tabular}{@{}p{0.45\linewidth}@{}p{0.45\linewidth}@{}>{\raggedleft\arraybackslash}p{0.10\linewidth}@{}}
    A) \ifnum#7=1 \textbf{#2} \else #2 \fi &
    B) \ifnum#7=2 \textbf{#3} \else #3 \fi & \\
    \multicolumn{2}{c}{} & \textbf{#6 pt} \\
    C) \ifnum#7=3 \textbf{#4} \else #4 \fi &
    D) \ifnum#7=4 \textbf{#5} \else #5 \fi & \\
  \end{tabular}
  \vspace{1.2em}
}

% QCM image (bonne réponse en gras)
\newcommand{\QCMimageCorr}[8]{%
  \stepcounter{qcmno}%
  \noindent
  \begin{minipage}[t]{0.69\linewidth}
    \textbf{\theqcmno.}~#1\par
  \end{minipage}%
  \hspace{0.02\linewidth}%
  \begin{minipage}[t]{0.29\linewidth}
    \vspace{-2pt}%
    \includegraphics[width=\linewidth]{#2}
  \end{minipage}

  \vspace{0.5em}
  \begin{tabular}{@{}p{0.45\linewidth}@{}p{0.45\linewidth}@{}>{\raggedleft\arraybackslash}p{0.10\linewidth}@{}}
    A) \ifnum#8=1 \textbf{#3} \else #3 \fi &
    B) \ifnum#8=2 \textbf{#4} \else #4 \fi & \\
    \multicolumn{2}{c}{} & \textbf{#7 pt} \\
    C) \ifnum#8=3 \textbf{#5} \else #5 \fi &
    D) \ifnum#8=4 \textbf{#6} \else #6 \fi & \\
  \end{tabular}
  \vspace{1.2em}
}

% ============================
% SECTION
% ============================

{% for item in qcms %}
  {% if item.has_image %}
    {% if is_correction %}
\QCMimageCorr{ {{ item.statement|tex_safe }} }
{ {{ item.image_path }} }
{ {{ item.A|tex_safe }} }
{ {{ item.B|tex_safe }} }
{ {{ item.C|tex_safe }} }
{ {{ item.D|tex_safe }} }
{ {{ item.points|default:0|floatformat:0 }} }
{ {{ item.correct_index|default:1 }} }
    {% else %}
\QCMimage{ {{ item.statement|tex_safe }} }
{ {{ item.image_path }} }
{ {{ item.A|tex_safe }} }
{ {{ item.B|tex_safe }} }
{ {{ item.C|tex_safe }} }
{ {{ item.D|tex_safe }} }
{ {{ item.points|default:0|floatformat:0 }} }
    {% endif %}
  {% else %}
    {% if is_correction %}
\QCMtexteCorr{ {{ item.statement|tex_safe }} }
{ {{ item.A|tex_safe }} }
{ {{ item.B|tex_safe }} }
{ {{ item.C|tex_safe }} }
{ {{ item.D|tex_safe }} }
{ {{ item.points|default:0|floatformat:0 }} }
{ {{ item.correct_index|default:1 }} }
    {% else %}
\QCMtexte{ {{ item.statement|tex_safe }} }
{ {{ item.A|tex_safe }} }
{ {{ item.B|tex_safe }} }
{ {{ item.C|tex_safe }} }
{ {{ item.D|tex_safe }} }
{ {{ item.points|default:0|floatformat:0 }} }
    {% endif %}
  {% endif %}
{% empty %}
\noindent\emph{Aucun QCM disponible.}
{% endfor %}
